% CS Labs Report --- Laboratory Work No. 3
\documentclass[12pt,a4paper]{article}

% Designed for XeLaTeX (Tectonic) and Overleaf.
\usepackage{fontspec}
\setmainfont{Liberation Serif}
\setmonofont{Liberation Mono}[Scale=MatchLowercase]
\usepackage[english]{babel}
\usepackage[left=2cm,right=1.5cm,top=2cm,bottom=2cm]{geometry}
\usepackage{setspace}
\usepackage{graphicx}
\usepackage{xcolor}
\usepackage{listings}
\usepackage{titlesec}
\usepackage{caption}
\usepackage{float}
\usepackage{amsmath}
\usepackage{amssymb}
\usepackage[hidelinks]{hyperref}

\graphicspath{{figures/}}
\onehalfspacing
\setlength{\parindent}{1.25cm}
\setlength{\parskip}{0pt}
\setlength{\emergencystretch}{2em}

\definecolor{codebackground}{HTML}{F3F4F6}
\definecolor{codeframe}{HTML}{D1D5DB}
\definecolor{commentgreen}{HTML}{3F6212}
\definecolor{keywordblue}{HTML}{1D4ED8}

\captionsetup[figure]{name=Fig.,labelsep=period,font=small}

\titleformat{\section}
  {\normalfont\bfseries\centering\fontsize{13}{16}\selectfont}
  {\thesection}{0.8em}{\MakeUppercase}
\titleformat{name=\section,numberless}
  {\normalfont\bfseries\centering\fontsize{13}{16}\selectfont}
  {}{0pt}{\MakeUppercase}
\titleformat{\subsection}
  {\normalfont\bfseries\fontsize{12}{15}\selectfont}
  {\thesubsection}{0.7em}{}
\titlespacing*{\section}{0pt}{0pt}{12pt}
\titlespacing*{\subsection}{0pt}{10pt}{6pt}

\lstdefinelanguage{Go}{
  keywords={package,import,const,var,func,for,range,if,else,return,make,map,string,int,bool,rune,true,false},
  morecomment=[l]{//},
  morestring=[b]",
  sensitive=true
}

\lstdefinestyle{go}{
  language=Go,
  basicstyle=\ttfamily\footnotesize,
  keywordstyle=\color{keywordblue}\bfseries,
  commentstyle=\color{commentgreen}\itshape,
  stringstyle=\color{red!55!black},
  backgroundcolor=\color{codebackground},
  frame=single,
  rulecolor=\color{codeframe},
  breaklines=true,
  columns=fullflexible,
  keepspaces=true,
  showstringspaces=false,
  extendedchars=true,
  numbers=left,
  numberstyle=\tiny\color{gray},
  xleftmargin=0.5cm,
  framexleftmargin=0.4cm,
  aboveskip=7pt,
  belowskip=7pt
}

\lstdefinestyle{python}{
  language=Python,
  basicstyle=\ttfamily\footnotesize,
  keywordstyle=\color{keywordblue}\bfseries,
  commentstyle=\color{commentgreen}\itshape,
  stringstyle=\color{red!55!black},
  backgroundcolor=\color{codebackground},
  frame=single,
  rulecolor=\color{codeframe},
  breaklines=true,
  columns=fullflexible,
  keepspaces=true,
  showstringspaces=false,
  extendedchars=true,
  numbers=left,
  numberstyle=\tiny\color{gray},
  xleftmargin=0.5cm,
  framexleftmargin=0.4cm,
  aboveskip=7pt,
  belowskip=7pt
}

\newcommand{\reportfigure}[4][\textwidth]{%
  \begin{figure}[H]
    \centering
    \includegraphics[width=#1]{#2}
    \caption{#3}
    \label{#4}
  \end{figure}
}

\begin{document}

\begin{titlepage}
  \thispagestyle{empty}
  \begin{center}
    \IfFileExists{figures/LOGO_UTM.jpg}{%
      \includegraphics[width=3.4cm]{figures/LOGO_UTM.jpg}\par
      \vspace{0.25cm}
    }{}

    {\bfseries\fontsize{14}{17}\selectfont
      MINISTRY OF EDUCATION, CULTURE AND RESEARCH\par
      OF THE REPUBLIC OF MOLDOVA\par
      Technical University of Moldova\par
      Faculty of Computers, Informatics and Microelectronics\par
      Department of Computer Science and Information Technologies\par
    }

    \vspace{1.0cm}
    {\bfseries\fontsize{16}{19}\selectfont Dinu Plesu, FAF-241\par}

    \vspace{1.0cm}
    {\bfseries\fontsize{36}{40}\selectfont Report\par}
    \vspace{0.25cm}
    {\itshape\fontsize{20}{24}\selectfont Laboratory Work No. 3\par}
    {\bfseries\itshape\fontsize{20}{24}\selectfont Cryptography and Security\par}

    \vspace{0.9cm}
    {\fontsize{16}{20}\selectfont Topic: Polyalphabetic ciphers --- the Vigen\`ere cipher\par}
    {\fontsize{15}{19}\selectfont over the Romanian alphabet (task 3.2)\par}
  \end{center}

  \vfill
  \begin{flushright}
    {\fontsize{15}{18}\selectfont Checked by:}\par
    \textbf{Maia Zaica}
  \end{flushright}

  \vfill
  \begin{center}
    {\bfseries\fontsize{14}{17}\selectfont Chi\c{s}in\u{a}u --- 2026}
  \end{center}
\end{titlepage}

\pagenumbering{arabic}
\setcounter{page}{2}
\tableofcontents
\newpage

\section{Task}

The students with an even number in the group register solve task~3.2: implement the Vigen\`ere cipher over the Romanian alphabet ($n=31$, codes $0\ldots30$ from Table~2). The program must encrypt and decrypt a text given interactively with a keyword key $k_2$ of at least 7 Romanian letters, display the key codes ``$k_i$'' and apply the formulas $c_i=(m_i+k_i)\bmod 31$ and $m_i=(c_i-k_i)\bmod 31$. Invalid keys and invalid characters must be reported clearly.

\section{Theoretical contents}

The Vigen\`ere cipher is polyalphabetic: instead of a single substitution it applies a sequence of Caesar shifts, cycling through a keyword. Each letter of the keyword supplies its own shift, taken from the alphabet encoding table. For an alphabet of $n$ letters and a keyword $K=k_1k_2\ldots k_\ell$ of length $\ell$ (with codes $k_i$), the $i$-th plaintext letter $m_i$ is encrypted as

\begin{align*}
  c_i &= (m_i+k_i)\bmod n, & k_i &= \operatorname{code}(K[i\bmod \ell]),\\
  m_i &= (c_i-k_i+n)\bmod n.
\end{align*}

For Romanian $n=31$. Because the same plain letter encrypts to different cipher letters depending on its position, the flat frequency profile of a monoalphabetic cipher disappears; the ciphertext frequency roughly flattens towards the language average, which defeats simple frequency analysis. The cipher is nevertheless not unconditionally secure: the key length $\ell$ can be discovered with the Kasiski examination or the index of coincidence, after which the text decomposes into $\ell$ interleaved monoalphabetic Caesar ciphers that are each broken by frequency.

\section{Implementation}

Task~3.2 was implemented in Go (file \texttt{task\_vigenere.go}). The keyword is validated (at least 7 Romanian letters), converted to codes over Table~2, and encryption/decryption loop over the repeating key:

\begin{lstlisting}[style=go,caption={Key codes and the Vigen\`ere loop (\texttt{task\_vigenere.go})}]
// key codes: each keyword letter -> its code in Table 2 (n = 31)
codes, ok, msg = vigenereKeyCodes(kw)
// encryption: ci = (mi + ki) mod 31 ; decryption: mi = (ci - ki) mod 31
for i, letter := range text {
    k := codes[i%len(codes)]
    if operation == opEncrypt {
        out[i] = alphabet[(letterCodes[letter]+k)%31]
    } else {
        out[i] = alphabet[(letterCodes[letter]-k+31)%31]
    }
}
\end{lstlisting}

The program prints the normalized text, the processed keyword, the key codes $k_i$, the formula used and the resulting ciphertext or plaintext, so every step can be checked by hand.

The same program also implements task~3.1, the Playfair cipher over the same 31-letter alphabet (file \texttt{task\_playfair.go}). Because 31 letters do not fit a classic $5\times5$ Playfair grid, the handout allows another shape; a $6\times6$ matrix is used, holding all 31 letters (J is kept as its own letter rather than folded onto I) plus the digits $0$--$4$ as fillers, for exactly 36 cells. The message is split into digraphs, a filler \texttt{X} is inserted between doubled letters and at an odd tail, and the rectangle / row / column rules of Playfair are applied. This task is a bonus --- the assigned variant for this student is 3.2 --- but it is fully implemented and verified in Fig.~\ref{fig:playfair}.

\section{Results of the work}

\reportfigure{fig1_encrypt.png}{Task 3.2 --- encryption with the keyword \texttt{ROMANIA}; key codes \texttt{[20 17 15 0 16 10 0]} and $c_i=(m_i+k_i)\bmod 31$.}{fig:encrypt}

\reportfigure{fig2_decrypt.png}{Task 3.2 --- decryption of the same ciphertext returns the original message.}{fig:decrypt}

\reportfigure{fig3_validation.png}{Validation --- a keyword shorter than 7 letters and a keyword containing digits are rejected.}{fig:validation}

\reportfigure{fig4_playfair_bonus.png}{Task 3.1 (bonus) --- the Playfair cipher over the same alphabet in a $6\times6$ matrix: the 31 letters plus the filler digits fill all 36 cells, and decryption restores the original message.}{fig:playfair}

The runs confirm both correctness and usability: Fig.~\ref{fig:encrypt} prints the processed keyword and its key codes [20 17 15 0 16 10 0] so that the shift applied at every position is transparent and hand-checkable, Fig.~\ref{fig:decrypt} proves the round trip is lossless, Fig.~\ref{fig:validation} shows the validation rules rejecting a short keyword and a keyword containing digits, each with an explicit message, and Fig.~\ref{fig:playfair} shows the bonus Playfair implementation (task 3.1) round-tripping a message through its $6\times6$ matrix.

\section{Conclusions}

The Vigen\`ere cipher over the Romanian alphabet was implemented and verified: encryption and decryption are exact inverses and all input constraints (keyword length, admissible characters) are enforced with explicit messages. Compared with the monoalphabetic ciphers of Labs~1--2, the polyalphabetic structure hides the letter frequencies --- the ciphertext frequency profile flattens --- so a plain frequency attack no longer works directly. However the cipher is breakable once the key period is found (Kasiski / index of coincidence), which motivates the one-time pad and modern stream ciphers. As a bonus, task~3.1 (Playfair over a $6\times6$ matrix) was also implemented and verified in the same program.

\section*{References}
\addcontentsline{toc}{section}{References}

\begin{enumerate}
  \item Lucrare de laborator nr.~3 --- Cifruri polialfabetice. UTM, FCIM --- handout (\texttt{resources/}).
  \item Table~2 --- letter encoding of the Romanian alphabet ($n=31$).
  \item W.~Stallings. \textit{Cryptography and Network Security} --- Vigen\`ere cipher and cryptanalysis.
\end{enumerate}

\section*{Source Code}
\addcontentsline{toc}{section}{Source Code}

The complete source code of all laboratory works (Go sources for Labs~1--3, Python sources for Labs~4--6, together with the ciphertext and plaintext files used in Lab~2 and Lab~6) is available in the GitHub repository:

\begin{center}
  \url{https://github.com/Prince-of-Nothing/CS-labs.git}
\end{center}

\end{document}
